\section{Related Work}
\label{sec:rel-work}
% \subsection{Learned Index}
% Learned indexes leverage the linearity of data distributions to map keys to possible positions, and this can be achieved using a non-linear model \cite{case18, cdfshop20, lipp21} or a piecewise linear model \cite{radixspline20, xindex20, madex20}. Currently, most learned indexes employ PLM to predict positions due to the simplicity of linear model parameters and the efficiency of training. In addition, some experiments have shown their effectiveness in various real-world datasets \cite{sosd20}. Moreover, learned indexes have also been extended to handle multimodal data, including one-dimensional key-value pairs \cite{pgm20, alex20, swix24, shift21}, high-dimensional vectors \cite{flood20, Tsunami20, coax23, learnedaknn22}, string data \cite{sindex20, rss21, lits24}, spatial data \cite{ifx20, lisa20, wisk23}, and more.

% \textbf{Segment Index:} In PLM-based learned indexes, the higher-level segment index plays a crucial role in overall performance. The RMI structure \cite{case18, pgm20} recursively constructs multilayer models by gradually magnifying the precision constraint to route to the underlying data. The tree-based hierarchical structure \cite{alex20, lipp21} divides and constructs the hierarchy more rapidly, and the lower coupling between layers allows local updates without training the entire model. In recent years, some learned indexes \cite{finedex21, swix24} have introduced flat structures to significantly reduce the depth of the routing by expanding the size of a single layer, but still require model training for efficient routing of elements within the layer. We propose the LM structure to address this issue. By computing the mapping parameters of elements, LM can locate any element within a flat layer with $O(1)$ cost without training a model, while also enabling efficient local updates.

% \textbf{Piecewise Linear Model}: Another major focus of current learned indexes is training PLM. PLM training can be divided into segmentation and regression. In segmentation, existing work either takes top-down approaches \cite{case18, alex20, swix24, carmi22} or bottom-up approaches \cite{fiting-tree19, pgm20, bourbon20}. The top-down method allows for fine-tuning the segments as a whole, such as SWIX \cite{swix24}, which uses second-order derivatives for segmentation to make the distribution within segments more linear. In contrast, the bottom-up approach segments the data in a data-driven manner based on certain constraints. PGM \cite{pgm20} introduces the \cite{line81} algorithm to ensure the minimum number of segments under a given error threshold constraint. In regression, \cite{alex20, sindex20, finedex21, swix24} employ the least squares method to minimize the average error, while \cite{pgm20, madex20, towards21} compute the linear parameters quickly by finding one line of the convex hull. LMIndex adopts a bottom-up approach to ensure the optimal number of segments and further utilizes OETA to efficiently obtain the linear parameters with the smallest maximum error within each segment.

% % \subsection{In-memory Index}
% % Tree-based in-memory indexes \cite{btree06, btree11, bwtree13} are mainstream techniques in traditional database systems. Index structures are organized in a tree-like manner, requiring significant memory and supporting both range and point queries. However, when indexing large data, as the tree depth increases and access paths grow, performance degradation can occur due to pointer chasing and cache misses. Hash-based in-memory indexes \cite{dynamichash98, hash16} excel in point query scenarios but lack efficient range query support, and hash collisions may lead to performance degradation. Hybrid index structures, such as LSM-tree \cite{lsmtree96, lsmtree17, remix21} and adaptive indexes \cite{adaptive18, adaptive22}, combine the advantages of various indexing techniques, adapting to different workloads. However, these structures are often more complex and have a narrower range of applicability.

\textbf{Evolution of Learned Indexes.}
The paradigm of learned indexing was pioneered by RMI \cite{case18}, which leverages the cumulative distribution function (CDF) of underlying data to map search keys to approximate physical locations. While early model-based indexes \cite{radixspline20, plex21, shift21} primarily targeted static workloads, the inherent limitations of read-only structures in dynamic environments spurred the development of updatable architectures. FITing-Tree \cite{fiting-tree19} and PGM \cite{pgm20} introduced bottom-up construction strategies governed by strict error bounds. Conversely, ALEX \cite{alex20} expanded upon the RMI concept by employing a top-down, data-driven construction that deliberately sacrifices structural balance to achieve superior model fitting across varying data granularities. These foundational works triggered a rapid proliferation of learned index designs, establishing a vast design space characterized by distinct structural layouts and optimization objectives.

\textbf{Structural Paradigms.}
The underlying routing architecture fundamentally dictates the search efficiency and update capability of a learned index. Hierarchical models typically employ either recursive, multi-layer error bounds \cite{case18, pgm20} or tree-based layouts \cite{alex20, lipp21} to iteratively refine position estimates. While tree-based designs facilitate localized updates without requiring holistic model retraining, the cascading traversal overhead remains non-trivial. Recent advancements \cite{uplif25, swix24, cabin24} advocate for flattened structures, deliberately reducing routing depth by expanding the capacity of individual layers. However, navigating these expansive layers still demands computational overhead for model inference. In contrast, our proposed LM structure addresses this bottleneck by computing deterministic mapping parameters. This enables $O(1)$ intra-layer routing without continuous model retraining, inherently supporting efficient local modifications while maintaining high accuracy.

\textbf{Model Construction and Optimization.}
Another critical dimension in learned index design is the methodology used for data segmentation and linear approximation. Segmentation strategies broadly bifurcate into top-down \cite{case18, alex20, swix24, carmi22} and bottom-up \cite{fiting-tree19, pgm20, bourbon20} frameworks. Top-down methodologies optimize segments holistically; for instance, SWIX \cite{swix24} utilizes second-order derivatives to enforce tighter linearity within segments. Conversely, bottom-up algorithms generate data-driven partitions governed by strict error thresholds. Within these segments, the regression techniques vary significantly: several frameworks optimize for average error using least squares \cite{alex20, sindex20, finedex21, swix24}, while others compute bounding parameters via convex hull algorithms to guarantee maximum error bounds \cite{pgm20, madex20, towards21}. Model retraining also incurs significant tail latency, a challenge that several studies \cite{chameleon24, vega25, hire25} have attempted to mitigate. Building upon these foundations, LMIndex leverages a bottom-up segmentation strategy coupled with the OETA algorithm to efficiently derive parameters that strictly minimize the maximum error per segment. Furthermore, this can be further refined through advanced parameter tuning \cite{tuning22, airindex23, newtuning25} and dynamic statistics maintenance \cite{sali23}.

\textbf{Expanding Modalities and Hardware Utilization.}
Beyond one-dimensional numeric keys, the learned index design space has diversified to accommodate complex data types. Researchers have tailored models for multi-dimensional vectors \cite{flood20, Tsunami20, coax23, learnedaknn22}, string data \cite{sindex20, rss21, lits24}, spatial datasets \cite{ifx20, lisa20, wisk23}, and streaming windows \cite{flirt23}. Concurrently, orthogonal optimizations have emerged to exploit modern hardware topologies, including CPU cache optimization \cite{carmi22}, GPU acceleration \cite{hyper24}, concurrent execution \cite{xindex20, finedex21}, and persistent memory efficiency \cite{apex21, simple23, learnedpmem23, disklearned24, wipe24, making24}. Despite these specialized advancements, balancing high-dimensional performance with robust, dynamic updatability remains an ongoing challenge in the field.
